\section{Conclusions}

A small scale turbopump, consisting of low specific speed Barske pump driven by a supersonic partial admission impulse turbine, was designed, built and tested, and its measured behaviour was used to answer the three objectives set in \cref{sec:intro_objectives}.

\textbf{Objective 1: Pump characterisation.} 

The Barske pump delivers its design head, with the measured head coefficient of \valPumpPsiMeas{} sitting 3\% above the design value and collapsing across five speeds, which demonstrates dynamic similarity. The peak measured efficiency is \valPumpEtaPeak, and the dominant loss was identified as disc frction, which absorbs about two thirds of the shaft power at test speed, at 2.5 to 2.9 times Barske's disc friction estimate, measured directly at shutoff. The cavitation behaviour was characterised and, through the transparent casing, visualised. The impeller eye cavitates early and benignly, while breakdown is driven by the diffuser throat, with the 3\% breakdown and the visual throat onset agreeing to around 0.1 m. Against the design methodology, Lock's model reproduces the measured head droop and locates the cavitation limited runout at the diffuser. However, the model predicts a higher effieicny than measured, which could be attributed either to mechanisms like tip leading edge rounding or other factors not accounted for in the model The affinity laws hold very well, and the fitted Reynolds scaling carries the efficiency to 26.4\% at the design speed. On the Barske architecture it delivers its design head, and its suction behaviour is dominated by a single identifiable mechanism that a transparent casing can watch. Whether around a quarter efficiency at design speed is acceptable depends entirely on the application, and for expendable small rocket stages where feed pressure mass dominates, it could be.

\textbf{Objective 2: Turbine design and model validation.} 

The turbine was designed with MoC blade generation, the Sasman and Cresci separation check, the Goldman starting criterion, and the Weiss loss models, giving a predicted total to static efficiency of \valTurbEtaTsDesign{} against an ideal 65\%. The standalone brake test was not performed, so the started stage efficiency remains a prediction. The starting criterion requires higher margins, as failure to start explained the coupled test result.

\textbf{Objective 3: Integration.} 

Coupled turbopump operation was demonstrated at around \valCoupledRpm{} RPM, with the turbine driving the pump to a measurable head rise and the shaft. The gap to the design speed is explained by two named and quantified mechanisms, the supply pressure droop to roughly 38\% of the design stator pressure and the predicted rotor unstart below approximately \valTurbStartRpm{} RPM. Taken together, the coupled result shows that the models used here is sufficient for design stage prediction of this class of machine. 


Overall, this work provides an open, experimentally validated characterisation of a Barske turbopump at $n_q \approx \valPumpNqDesign$, including a measured loss breakdown and a direct visualisation of the cavitation mechanism that sets its operating limit.


\subsection{Future Work}

% \todo{future work: find the cavitation curve that lock describes, which can only be seen by slowly closing the outlet valve instead of what I did which is slowly closing inlet valve. basically, just more test as none of my data was able to capture what he found... unless I need to analyze the H-Q curves. Yes I will do that.}

\begin{enumerate}
    \item Test a turbine designed with a starting inlet mach number lower than the goldman starting criterion to provide results that can properly compared to the models. A started turbine would resolve the single largest open prediction, the started stage efficiency of \valTurbEtaTsDesign.
    % \item A 15 mm direct rigid coupler to unlock the full 20,000 RPM envelope of the torquemeter, removing the drive side speed limit on pump testing.
    \item An enlarged diffuser throat, traded against suction margin. The measured throat limited runout and the published optimisation trend \cite{danieliConceptualDesignElectric2023} both point at the same lever.
    \item Hot gas generator products as the working fluid, closing the specific heat ratio gap between cold air at $\gamma = 1.4$ and combustion gas near 1.2, which also shifts the starting margin.
    % \item A machined metal impeller, to isolate the contribution of the printed surface quality to the measured churning multiplier.
    % \item A full admission variant at larger diameter, to isolate the partial admission sector losses from the rest of the loss budget.
    \item CFD of the pump geometry with the measured boundary conditions, now that there is experimental data to validate it against.
    \item CFD of the turbine geometry to gain a better understanding of the low measured efficiency and the unstart mechanism, and to provide a more detailed loss breakdown than the Weiss model.
    % \item Global shutter imaging with a synchronised strobe through the existing transparent casing. The present consumer camera resolves the cavity envelope but not bubble dynamics, since the frame rate sits far below blade passing and the rolling shutter distorts fast structures; individual bubbles already appear in isolated frames, so the optical access is proven and only the camera limits what it can measure.
    % \item High rate burst logging on the pressure channels. The logged data are 32 sample averages at 50 Hz, so blade passing unsteadiness at several hundred hertz is unresolvable by design; short bursts at the ADC's native rate would resolve it without changing the architecture.
\end{enumerate}